\section{Gossip dissemination}\label{sec:session_gossip}

Nodes keep their local CRDT states (section~\ref{sec:session_crdt}) aligned with
each other by exchanging \emph{gossip} messages.
A gossip carries a single topic record and is sent over one of four carriers
described below.

\subsection{Gossip carriers}

\begin{description}
    \item[Broadcast gossip] Multicast on the broadcast subject
    (section~\ref{sec:session_subject_id_mapping}). Every node is expected to be
    attached to this subject, so a broadcast gossip is eventually visible to every
    participant; broadcast traffic is the primary mechanism for escaping
    partitioned or newly-joined subsets. The broadcast subject is also used for
    scout messages (section~\ref{sec:session_scout}).

    \item[Sharded gossip] Multicast on a \emph{gossip shard} subject-ID.
    A topic is assigned to the shard
    \[
        \mathit{shard\_index} = \mathit{hash} \bmod \mathit{gossip\_shard\_count}
    \]
    where $\mathit{gossip\_shard\_count}$ is the number of subject-IDs between
    the auto-allocated range and the broadcast subject (see
    section~\ref{sec:session_subject_id_mapping}). A node joins every gossip shard
    that any of its local topics maps into, and no others.

    \item[Inline gossip] A gossip state carried in the header of a
    \texttt{msg\_be} or \texttt{msg\_rel} application message
    (section~\ref{sec:session_headers}). Inline gossip is observed by every
    subscriber of the topic and converges the CRDT state for that topic without
    spending a separate message on it. Handling of inline gossip is recommended
    but not required for correctness.

    \item[Unicast gossip] A single gossip sent to a specific remote node,
    used for scout responses and for targeted repair of a known divergence.
\end{description}

\subsection{Schedules}

A node that holds at least one explicit topic emits gossip on three schedules:
periodic, urgent, and inline. Implicit topics
(section~\ref{sec:session_patterns}) do not emit gossip.

\paragraph{Periodic gossip} Every explicit topic emits its current state at a
randomized interval around a nominal period (recommended default: $5$ seconds,
randomized by $\pm$ one eighth).
Every $N$-th emission goes on the broadcast subject rather than the topic's
gossip shard, so the broadcast lane sees a slow but guaranteed sweep of all
topics on the network (recommended $N = 10$).
Between broadcast sweeps the topic's shard is used, which scales better on large
networks because each node only joins a fraction of the shards.

\paragraph{Urgent gossip} Whenever a topic's local subject-ID occupancy changes
-- a new topic is created, an incoming gossip forces a re-allocation, or a
cascade evicts a record -- a broadcast gossip is scheduled with a small random
delay $W_u$ (section~\ref{sec:session_gossip_delays}). If an equivalent or
dominating gossip from another node arrives before the timer fires, the pending
urgent gossip is suppressed.

\paragraph{Inline gossip} A node that publishes a \texttt{msg\_be} or
\texttt{msg\_rel} message piggybacks the current gossip state of the message's
topic in the header. A receiver that elects to process inline gossip applies it
to its local CRDT in the usual way; this is the fastest possible CRDT update on
active topics because it consumes no additional network capacity.

\subsection{Suppression rules}\label{sec:session_gossip_suppression}

Gossip schedules are designed so that the network-wide rate of gossip traffic on
any given topic stays roughly constant regardless of the number of nodes holding
it. To that end, nodes apply two suppression rules when incoming gossip indicates
that some other node is already covering the same information:

\begin{itemize}
    \item \textbf{Cancel a pending urgent gossip} when an incoming gossip arrives
    on the broadcast subject whose arbitration priority dominates the state the
    urgent gossip was intended to publish. The incoming broadcast already
    carries the information forward, so repeating it would waste network
    capacity.

    \item \textbf{Slow down periodic gossip} when an incoming non-inline gossip
    matches the local $\mathit{lage}$ and eviction counter for the same topic.
    Another node is already covering the topic at the current cadence; the
    local node can afford to reschedule its next emission with a longer
    interval.
\end{itemize}

\begin{remark}
    Inline gossip is deliberately excluded from the periodic slowdown rule.
    Because inline gossip is piggybacked on application messages, its rate is
    determined by the application's publication rate, not by the gossip
    scheduler; allowing it to suppress periodic gossip would couple the CRDT
    repair rate to the application traffic pattern.
\end{remark}

\subsection{Randomized delay windows}\label{sec:session_gossip_delays}

The urgent window $W_u$ and the startup window $W_s$ are both tunable. The
recommended orders of magnitude below are derived from the analytic model in
\texttt{model/gossip\_propagation.ipynb} of the reference implementation.

Let $\delta_{\max}$ be the worst-case propagation and processing delay for a
gossip between two nodes on the network.

\begin{description}
    \item[Urgent window $W_u$] Controls the amount of jitter inserted before an
    urgent gossip. If $m_u$ nodes independently schedule an urgent repair for
    the same event, the expected number of redundant urgent emissions is
    approximately $(m_u - 1) \cdot \delta_{\max} / W_u$.
    In typical deployments $m_u$ is small because a given collision or
    divergence is observed by only a few nodes; setting
    $W_u \approx \delta_{\max}$ is usually sufficient.

    \item[Startup window $W_s$] Controls the amount of jitter inserted when a
    node comes online and schedules its first round of gossip.
    Wider windows reduce the probability of simultaneous emissions from many
    startups but saturate quickly; recommended range is
    $[\delta_{\max},\ 10 \cdot \delta_{\max}]$.
\end{description}

On a typical local network with $\delta_{\max}$ on the order of tens of
milliseconds, choosing $W_u$ on the order of tens of milliseconds and $W_s$ on
the order of tens to a few hundred milliseconds is sufficient.

\subsection{Implicit topic retirement}

An \emph{implicit} topic is one that a node holds only because a local pattern
subscription (section~\ref{sec:session_patterns}) matched a remote gossip. An
implicit topic does not emit gossip, and it is automatically retired once it has
observed no messages and no gossip activity for the implicit-topic retention
window. The recommended retention window is ten minutes; it should be much
longer than the worst-case periodic gossip interval so that transient silence
on an active topic does not cause premature retirement.

\section{Scout}\label{sec:session_scout}

The scout mechanism lets a node pull in topics that match a pattern without
waiting for periodic gossip rotation. A scout is a session-layer message whose
payload is a pattern string (section~\ref{sec:session_names}).

A scout message is broadcast on the broadcast subject using the
\texttt{scout} header (section~\ref{sec:session_headers}).
Every node that receives the message checks its local topics against the
pattern; for each match, the receiver replies with a \texttt{gossip} header
describing the matched topic. Replies are typically unicast back to the
originating node, using the unicast context provided by the transport with the
incoming scout, but broadcasting is also permitted if that is easier for the
responder.

A scout with an empty pattern is reserved for a future extension -- a request
for a network-wide health summary -- and shall be discarded by current
implementations.
